\documentclass{beamer}

\usepackage{polyglossia}

\usepackage[orientation = portrait, size = a1, scale = 1.4]{beamerposter}
%\usepackage[backend = biber, style = iso-authoryear, sortlocale = en_US, autolang = other, bibencoding = UTF8]{biblatex}
\usepackage{datetime}
\usepackage{adjustbox}
%\usepackage[utf8]{inputenc}
\usepackage{fontspec}
\usepackage{mathrsfs}
\usepackage{microtype}
\usepackage{listings}
\usepackage{booktabs}

\setdefaultlanguage{english}
\setmainfont{TeX Gyre Termes}
\usetheme{gemini}
\usecolortheme{mit}

\usepackage{tikz}
\usepackage{pgfplots}
\usetikzlibrary{arrows,automata,shapes,calc, patterns, backgrounds}

\definecolor{gray}{rgb}{0.5,0.5,0.5}
\definecolor{mauve}{rgb}{0.58,0,0.82}
\lstset{basicstyle={\small\ttfamily},
	belowskip=3mm,
	breakatwhitespace=true,
	breaklines=true,
	classoffset=0,
	columns=flexible,
	framexleftmargin=0.25em,
	frameshape={}{}{}{}, %To remove to vertical lines on left, set `frameshape={}{}{}{}`
	keywordstyle=\color{blue},
	numbers=none, %If you want line numbers, set `numbers=left`
	numberstyle=\color{gray},
	showstringspaces=false,
	stringstyle=\color{mauve},
	tabsize=3,
	xleftmargin =1em
}

\DeclareMathOperator*{\argmax}{arg\,max}
\def\mathdefault#1{#1}

\title{Benchmarking and Transfer Learning for Hyperparameter Optimization of Graph Neural Networks}

\author{%
	Marek Dědič\inst{1 2} \and
	Michal Bělohlávek\inst{1 2}
}
\institute{%
	\inst{1} Czech Technical University in Prague \samelineand
	\inst{2} Cisco Systems, Inc.
}

\footercontent{%
	COSEAL 2026, Aachen, Germany \hfill
	\href{mailto:marek@dedic.eu}{marek@dedic.eu}
}

% \logoright{\includegraphics[height=7cm]{logo1.pdf}}
% \logoleft{\includegraphics[height=7cm]{logo2.pdf}}

% If you have N columns, choose \sepwidth and \colwidth such that
% (N+1)*\sepwidth + N*\colwidth = \paperwidth
\newlength{\sepwidth}
\newlength{\colwidth}
\setlength{\sepwidth}{0.033\paperwidth}
\setlength{\colwidth}{0.45\paperwidth}
\newcommand{\separatorcolumn}{\begin{column}{\sepwidth}\end{column}}

\newcommand{\corr}{(\Letter)}
\newcommand{\name}[1]{\textit{#1}}
\newcommand{\mathfield}{\ensuremath{\mathbb}}
\newcommand{\mathmat}{\ensuremath{\mathbf}}
\newcommand{\mathset}{\ensuremath{\mathbb}}
\newcommand{\mathspace}{\ensuremath{\mathcal}}
\newcommand{\mathvec}{\ensuremath{\bm}}

\begin{document}
\begin{frame}[fragile,t]

\begin{columns}[t]
	\separatorcolumn

	\begin{column}{\colwidth}
		\begin{alertblock}{3 main projects}
			\begin{itemize}
				\item A HARP-based method for solving performance-complexity trade-off in GNNs via adaptive prolongation.
				\item CSP is a very simple learning method for classification and retrieval on graphs, hyper-graphs or problems described by one-hot features. It achieves comparable performance to reference methods in multiple problems, while being significantly cheaper to compute.
				\item A hyper-parameter optimizer benchmark for GNNs and a novel Cross-RF transfer metalearning method that outperforms SMBOs.
			\end{itemize}
		\end{alertblock}
	\end{column}

	\separatorcolumn

	\begin{column}{\colwidth}
		\begin{block}{Hyperparameter optimization for GNNs}
			We benchmarked 5 standard HPO methods on 9 GNN datasets. The results show that Bayesian optimization (BO) and Tree-structured Parzen Estimators (TPE) achieve the best performance, while Quasi Monte-Carlo (QMC) falls behind.
			\adjustbox{width=\columnwidth}{%
				\input{images/benchmark_ranks/benchmark_ranks.pgf}
			}
		\end{block}

		\begin{block}{The Cross-RF method for HPO on graphs}
			We propose a HPO method using a metamodel \( \mathcal{M}_\rho : \mathfield{R}^d \times \boldsymbol\Lambda \to \mathfield{R} \) that takes as input the properties of a graph dataset \( \delta \left( \mathcal{D} \right) \) and a hyperparameter configuration \( \lambda \), and outputs an estimate of a performance metric \( \rho \). This gives us the following HPO procedure:
			\begin{equation*}
				\tau \left( \mathcal{D}, \mathscr{F}, \tilde{\Lambda}, \rho \right) = \argmax_{\lambda \in \tilde{\Lambda}} \mathcal{M}_\rho \left( \delta \left( \mathcal{D} \right), \lambda \right)
			\end{equation*}
			We show this method to outperform all of the above methods on the 9 graph datasets by a slight margin.
		\end{block}
	\end{column}

	\separatorcolumn
\end{columns}

\end{frame}
\end{document}
